% Generated by task tex-libraries from dossiers/lib-others.md, section B (and F.2, F.3 for the
% new pin). Snippets are copied from the dossier, which copied them from the sources.
\begingroup\sloppy
\section{CompPoly: tables, polynomials and the two fields}\label{sec:gen-lib-comppoly}

CompPoly is the library of \emph{computable} polynomials: representations that run
(\texttt{\#eval}, compiled code) and are proved equal to Mathlib's non-computable ones. Every snippet
of this section is from CompPoly at the old pin \code{3468b38c} (\enquote{feat(fields): add
polynomial-basis GF(2\textasciicircum 64) and its cubic extension (\#321)}), read with
\code{git show 3468b38c:<path>}; \cref{sec:gen-lib-comppoly-new} says what the new pin
\code{572f9973} changes.

\subsection{A multilinear polynomial as its table of values}\label{sec:gen-lib-comppoly-mle}

\begin{define}{\lean{CMlPolynomialEval}: a table on the Boolean cube}
\lean{CMlPolynomialEval R n} is a table of $2^n$ values in \lean{R}; entry $i$ is the value at the
cube point whose coordinate $j$ is bit $j$ of $i$ (\textbf{bit 0 is variable 0}, little-endian).
It stands for the multilinear polynomial with those values on $\{0,1\}^n$, defined through the
Lagrange (equality) basis. Being \lean{@[reducible]}, it \emph{is} a \lean{Vector} to instance
search, which is why leanerVM wraps it in a structure \lean{Column}
(\cref{sec:gen-l0-oracle}).
\end{define}

\begin{leancode}
/-- `CMlPolynomialEval n R` is the type of multilinear polynomials in `n` variables over a ring `R`.
  It is represented by its evaluations over the Boolean hypercube `{0,1}^n`,
  i.e. Lagrange basis coefficients.
  The indexing is **little-endian** (i.e. the least significant bit is the first bit). -/
@[reducible]
def CMlPolynomialEval (R : Type*) (n : ℕ) := Vector R (2 ^ n) -- coefficient of Lagrange basis
\end{leancode}

\src{\file{CompPoly/}\allowbreak\file{Multilinear/}\allowbreak\file{Basic.lean:42-48}, CompPoly \code{3468b38c}}


\begin{leancode}
/-- Lagrange (hypercube) basis at point `w`.

Returns the length-`2^n` vector `v` such that for any `x ∈ {0,1}^n`, letting
`i = ∑_{j=0}^{n-1} x_j · 2^j` (little‑endian indexing), we have
`v[i] = ∏_{j < n} (x_j · w[j] + (1 - x_j) · (1 - w[j]))`.
Equivalently, for `i : Fin (2^n)`,
`v[i] = ∏_{j < n}, (if the j-th bit of i is 1 then w[j] else 1 - w[j])`.
-/
def lagrangeBasis (w : Vector R n) : Vector R (2 ^ n) :=
  Vector.ofFn (fun i => ∏ j : Fin n, if (BitVec.ofFin i).getLsb j then w[j] else 1 - w[j])
\end{leancode}

\src{\file{CompPoly/}\allowbreak\file{Multilinear/}\allowbreak\file{Basic.lean:401-411}, CompPoly \code{3468b38c}}


\begin{leancode}
/-- Evaluate a `CMlPolynomialEval` at a point -/
def eval (p : CMlPolynomialEval R n) (x : Vector R n) : R :=
  Vector.dotProduct p (lagrangeBasis x)
...
/-- Evaluate a `CMlPolynomialEval` at a point using a ring homomorphism -/
def eval₂ (p : CMlPolynomialEval R n) (f : R →+* S) (x : Vector S n) : S := eval (map f p) x

/-- Evaluate the multilinear equality kernel `eq̃(w, x)`. -/
@[inline] def eqTilde (w x : Vector R n) : R :=
  eval (lagrangeBasis w) x
\end{leancode}

\src{\file{CompPoly/}\allowbreak\file{Multilinear/}\allowbreak\file{Basic.lean:523-544}, CompPoly \code{3468b38c}}

So \lean{eval p x} $= \sum_i p[i] \cdot \prod_j (\mathrm{bit}_j(i)\ ?\ x_j : 1 - x_j)$, the
multilinear extension $\mle{p}(x)$, and \lean{eqTilde w x}
$= \sum_i \eq(w,i)\,\eq(x,i) = \prod_j \bigl(w_j x_j + (1-w_j)(1-x_j)\bigr)$. The fast evaluator
folds one variable at a time (variable 0 first; \lean{evalMleStep}, \file{:467-471}: entry $j$ of
the folded table is $(1 - x_0)\cdot\mathit{values}[2j] + x_0\cdot\mathit{values}[2j+1]$) and is
proved equal to the dot product. \lean{map f p} applies a ring homomorphism \lean{f : R →+* S} to
every entry (\file{:461-463}); with \lean{f = algebraMap K E} it lifts a $\K$-table to an
$\E$-table, so \lean{eval₂Mle q (algebraMap K E) r} evaluates the $\K$-multilinear $\mle{q}$ at
a point $r \in \E^n$.

\begin{leancode}
/-- Evaluate a `CMlPolynomialEval` by recursive multilinear-extension interpolation. -/
@[inline, specialize]
def evalMle (p : CMlPolynomialEval R n) (x : Vector R n) : R :=
  evalMleValues p x
...
/-- Evaluate a `CMlPolynomialEval` through a ring homomorphism using multilinear-extension
interpolation. -/
@[inline, specialize]
def eval₂Mle (p : CMlPolynomialEval R n) (f : R →+* S) (x : Vector S n) : S :=
  evalMle (map f p) x
\end{leancode}

\src{\file{CompPoly/}\allowbreak\file{Multilinear/}\allowbreak\file{Basic.lean:497-521}, CompPoly \code{3468b38c}}


\begin{leancode}
theorem eval_mle_eq_eval (p : CMlPolynomialEval R n) (x : Vector R n) :
    evalMle p x = eval p x := by
...
theorem eval₂_mle_eq_eval₂ (p : CMlPolynomialEval R n) (f : R →+* S) (x : Vector S n) :
    eval₂Mle p f x = eval₂ p f x := by
\end{leancode}

\src{\file{CompPoly/}\allowbreak\file{Multilinear/}\allowbreak\file{Basic.lean:585-586} and \file{:638-640}, CompPoly \code{3468b38c}}

\paragraph{Bridge to Mathlib.} At the CompPoly pin there is a conversion
\lean{CMlPolynomialEval.}\allowbreak\lean{toMvPolynomial} (\file{Multilinear/}\allowbreak\file{Equiv.lean:337}) but \textbf{no theorem
that \lean{eval} is Mathlib's evaluation of it}; that bridge is ArkLib's. The blueprint cites it
correctly as ArkLib's (\file{protocol-blueprint.md:246}). For a reader the definition \lean{eval}
above is self-explanatory and needs no bridge to be trusted.

\begin{leancode}
theorem eval_eq_MvPolynomial_MLE (evals : (Fin n → Fin 2) → R) (x : Fin n → R) :
    eval
        (Vector.ofFn fun i => evals (finFunctionFinEquiv.symm i))
        (Vector.ofFn x) =
      MvPolynomial.eval x (MvPolynomial.MLE evals) := by
\end{leancode}

\src{\file{ArkLib/}\allowbreak\file{ToCompPoly/}\allowbreak\file{Multilinear/}\allowbreak\file{Basic.lean:56}, ArkLib \code{dca90385}}

\paragraph{Bit order against leanVM.} The specification fixes \enquote{$\langle u\rangle :=
\sum_i u_i 2^i$ \dots\ (bit $0$ first)} (\file{doc/}\allowbreak\file{leanvm/}\allowbreak\file{body/}\allowbreak\file{b-polynomial-commitment-scheme.tex:278})
and places each table's selector in the \emph{high} bits of its offset
(\file{04-committing-the-witness.tex:10}, \enquote{$\mathit{sel}_i \in \{0,1\}^{M-\kappa_i}$ for
the high bits of that offset}). CompPoly's little-endian order agrees: the low variables of the
stack index the position inside a block.

\subsection{Computable multivariate polynomials}\label{sec:gen-lib-comppoly-mv}

\begin{define}{\lean{CMvPolynomial}: sparse multivariate polynomials that run}
\lean{CMvPolynomial n R} is a finite map from monomials (exponent vectors \lean{Vector ℕ n}) to
coefficients, with no zero coefficient stored (\lean{Unlawful n R} is the map, \lean{Lawful} the
subtype without zeros). The spine states its constraint and flush polynomials in this type
(\file{LeanerVM/}\allowbreak\file{Protocol/}\allowbreak\file{Spine/}\allowbreak\file{Instance.lean:123-132}, \file{Seams.lean:77}). \lean{eval₂},
\lean{eval} and \lean{totalDegree} are computed by folding over the map; a bijection with
Mathlib's \lean{MvPolynomial} (the standard, non-computable multivariate polynomial type)
transfers evaluation and degree.
\end{define}

\begin{leancode}
/-- A computable multivariate polynomial in `n` variables with coefficients in `R`. -/
abbrev CMvPolynomial (n : ℕ) (R : Type*) [Zero R] : Type _ := Lawful n R
...
/-- The subtype of polynomials with no zero coefficients. -/
@[implicit_reducible]
def Lawful (n : ℕ) (R : Type*) [Zero R] : Type _ :=
  {p : Unlawful n R // p.isNoZeroCoef}
\end{leancode}

\src{\file{CompPoly/}\allowbreak\file{Multivariate/}\allowbreak\file{CMvPolynomial.lean:47-48} and \file{Multivariate/}\allowbreak\file{Lawful.lean:33-36}, CompPoly \code{3468b38c}}


\begin{leancode}
def eval₂ {R S : Type*} {n : ℕ} [Semiring R] [CommSemiring S] :
    (R →+* S) → (Fin n → S) → CMvPolynomial n R → S :=
  fun f vs p => ExtTreeMap.foldl (fun s m c => (f c * MonoR.evalMonomial vs m) + s) 0 p.1
...
def eval {R : Type*} {n : ℕ} [CommSemiring R] : (Fin n → R) → CMvPolynomial n R → R :=
  eval₂ (RingHom.id _)
...
def totalDegree {R : Type*} {n : ℕ} [Zero R] : CMvPolynomial n R → ℕ :=
  fun p => Finset.sup (List.toFinset (List.map CMvMonomial.toFinsupp (Lawful.monomials p)))
    (fun s => Finsupp.sum s (fun _ e => e))
\end{leancode}

\src{\file{CMvPolynomial.lean:131-135}, \file{:216-218}, \file{:229-233}, CompPoly \code{3468b38c}}


\begin{leancode}
def fromCMvPolynomial  (p : CMvPolynomial n R) : MvPolynomial (Fin n) R :=
  let support : List (Fin n →₀ ℕ) := p.monomials.map CMvMonomial.toFinsupp
  let toFun (f : Fin n →₀ ℕ) : R := p[CMvMonomial.ofFinsupp f]?.getD 0
  let mem_support_fun {a : Fin n →₀ ℕ} : a ∈ support ↔ toFun a ≠ 0 := by grind
  AddMonoidAlgebra.ofCoeff <| Finsupp.mk support.toFinset toFun (by simp [mem_support_fun])
\end{leancode}

\src{\file{CompPoly/}\allowbreak\file{Multivariate/}\allowbreak\file{MvPolyEquiv/}\allowbreak\file{Core.lean:35-39}, CompPoly \code{3468b38c}}


\begin{leancode}
lemma eval_equiv {p : CMvPolynomial n R} {vals : Fin n → R} :
    p.eval vals = (fromCMvPolynomial p).eval vals := by
...
lemma totalDegree_equiv {S : Type*} {p : CMvPolynomial n R} [CommSemiring S] :
    p.totalDegree = (fromCMvPolynomial p).totalDegree := by rfl
\end{leancode}

\src{\file{CompPoly/}\allowbreak\file{Multivariate/}\allowbreak\file{MvPolyEquiv/}\allowbreak\file{Eval.lean:53-61}, CompPoly \code{3468b38c}}

With \lean{polyEquiv : CMvPolynomial n R ≃ MvPolynomial (Fin n) R} (\file{Core.lean:134}), a
degree bound \lean{C.totalDegree ≤ d} in the spine is Mathlib's total degree, and
Schwartz--Zippel from Mathlib applies through \lean{eval_equiv}. Building polynomials (\lean{C},
\lean{X}, ring operations) needs \lean{[BEq R] [LawfulBEq R]} (\file{CMvPolynomial.lean:55-62}).
$\K$ has it (instance search finds \lean{instLawfulBEq} from \lean{DecidableEq}; probe
\code{FieldFidelity}, last lines); \textbf{$\E$ has none at the old pin} (the same probe:
\code{failed to synthesize LawfulBEq E}), so the spine cannot build a \lean{CMvPolynomial} over
$\E$. The blueprint records this (\file{protocol-blueprint.md:570}) and keeps every polynomial
over $\K$, which matches leanVM (constraints have $\K$ coefficients).

\subsection{The fields \texorpdfstring{$\K = \mathrm{GF}(2^{64})$}{K} and \texorpdfstring{$\E = \mathrm{GF}(2^{192})$}{E}}\label{sec:gen-lib-comppoly-fields}

\begin{define}{\lean{BF64}: the field $\K$}
$\K = \mathrm{GF}(2)[x]/(x^{64} + x^4 + x^3 + x + 1)$. CompPoly keeps the specification object
and the computable carrier apart and joins them by a proved bijective ring homomorphism. The
specification is Mathlib's quotient \lean{AdjoinRoot basePoly}, a field because the modulus is
irreducible; the carrier is a 64-bit word whose bit $i$ is the coefficient of $x^i$, with
addition XOR and multiplication the carry-less product reduced modulo the modulus.
\end{define}

\begin{leancode}
/-- The modulus `x^64 + x^4 + x^3 + x + 1` over `GF(2)`. Part of the specification only;
it is never evaluated. -/
noncomputable def basePoly : Polynomial (ZMod 2) := X ^ 64 + X ^ 4 + X ^ 3 + X + 1
...
theorem basePoly_irreducible : Irreducible basePoly := by
  refine Polynomial.irreducible_of_rabin (d := 64) ?_ (by norm_num) ?_ ?_
...
noncomputable abbrev BF64Quot : Type := AdjoinRoot basePoly
\end{leancode}

\src{\file{CompPoly/}\allowbreak\file{Fields/}\allowbreak\file{Binary/}\allowbreak\file{BF64/}\allowbreak\file{Basic.lean:100-102}, \file{:152-153}, \file{:199}, CompPoly \code{3468b38c}}

Irreducibility is proved by Rabin's test (CompPoly's \lean{Polynomial.}\allowbreak\lean{irreducible_}\allowbreak\lean{of_rabin})
against certificate chains that the kernel checks by \lean{rfl} (\file{BF64/}\allowbreak\file{BaseCertificate.lean},
552 lines); no \lean{native_decide} (axioms of \lean{BF64.basePoly_irreducible}: \lean{propext},
\lean{Classical.choice}, \lean{Quot.sound}, probe \code{Layer0}).

\begin{leancode}
/-- `GF(2^64)` in its computable, machine representation: a 64-bit word whose bit `i` is
the coefficient of `x^i`. -/
abbrev BF64 : Type := BitVec 64

namespace BF64

instance : Zero BF64 := ⟨(0 : BitVec 64)⟩
instance : One BF64 := ⟨(1 : BitVec 64)⟩

/-- Addition in characteristic two is `xor`. -/
instance : Add BF64 := ⟨fun a b => a ^^^ b⟩
...
/-- Multiplication: the carry-less product, reduced modulo the modulus. -/
instance : Mul BF64 :=
  ⟨fun a b => reduce (carryLessMul (w := 128) a b)⟩
\end{leancode}

\src{\file{CompPoly/}\allowbreak\file{Fields/}\allowbreak\file{Binary/}\allowbreak\file{BF64/}\allowbreak\file{Impl.lean:52-71}, CompPoly \code{3468b38c}}

The carrier maps to the quotient by \lean{toQuot a = AdjoinRoot.}\allowbreak\lean{mk basePoly (toPoly a)}
(\file{:77-79}), proved a ring homomorphism (\lean{toQuot_add}, \file{:103}; \lean{toQuot_mul},
\file{:107}), injective (\file{:118}) and surjective (\file{:398}). Every ring and field law of
\lean{BF64} is transported through it (\lean{CommRing}, \file{:237-255}; \lean{Field},
\file{:442-450}); inversion is the Itoh--Tsujii chain, proved correct (\lean{mul_invItohTsujii},
\file{:404}), with $0^{-1} = 0$. \lean{DecidableEq K} is \lean{BitVec}'s
(\lean{instDecidableEqBitVec}, probe \code{Layer0}). Cardinality and the instances:

\begin{leancode}
/-- The carrier is in bijection with `Fin (2 ^ 64)`, by its underlying representation. -/
def equivFin : BF64 ≃ Fin (2 ^ 64) where
  toFun a := a.toFin
  invFun i := BitVec.ofFin i
  left_inv _ := rfl
  right_inv _ := rfl

instance : Fintype BF64 := Fintype.ofEquiv _ equivFin.symm

/-- `BF64` has `2 ^ 64` elements. -/
theorem card_bf64 : Fintype.card BF64 = 2 ^ 64 := by
  rw [Fintype.card_congr equivFin, Fintype.card_fin]
\end{leancode}

\src{\file{CompPoly/}\allowbreak\file{Fields/}\allowbreak\file{Binary/}\allowbreak\file{BF64/}\allowbreak\file{Impl.lean:384-395}, CompPoly \code{3468b38c}}

\paragraph{Hazard of the eager \texorpdfstring{\texttt{Fintype BF64}}{Fintype BF64} (old pin).}
The instance is a closed computable term; compiled code initialises it at program start by
building \lean{List.finRange (2^64)}, so any executable linking the module is killed for lack of
memory. leanerVM records this (\file{docs/}\allowbreak\file{roadmap/}\allowbreak\file{leanisa-status.md:753-769} at \code{b435631},
the finding \enquote{\lean{instance : Fintype BF64} is evaluated at executable startup} (P3)):
the test driver had to become a library and no leanerVM executable may link the field. It cannot
simply be made \lean{noncomputable} at this pin, because \lean{ExtensionParams F} takes
\lean{[Fintype F]} and so every \lean{Ext} operation carries it at run time. Consequence for the
blueprint: the executable verifier \lean{verify}, differential tests against the Rust, and any
native target are \textbf{blocked at this pin} until the CompPoly fix lands; it has
(\cref{sec:gen-lib-comppoly-new}).

\paragraph{Hazard of \texorpdfstring{\texttt{abbrev BF64 := BitVec 64}}{abbrev BF64} (old pin).}
Because \lean{BF64} is reducible, CompPoly's \lean{Add}, \lean{Mul}, \lean{Neg}, \lean{Inv},
\lean{NatCast}, \lean{Fintype}, \lean{CommRing}, \lean{Field} instances are instances \textbf{on
\lean{BitVec 64} itself} in every file that imports the field. The probe \code{FieldFidelity}
(lines 85--100, output lines 2--12) shows it:

\begin{leancode}
#eval u + v                -- u v : BitVec 64 := 3, 5   ⟹  6#64   (XOR, not 8)
#eval (u + v == 8, u + v == 6)                          ⟹  (false, true)
#synth Add (BitVec 64)                                  ⟹  BF64.instAdd
#synth Add (BitVec 32)                                  ⟹  BitVec.instAdd
\end{leancode}

\src{probe \code{FieldFidelity} (\file{probes/}\allowbreak\file{lib-others/}\allowbreak\file{FieldFidelity.lean}), lines 85--100 and output lines 2--12, run at the old pins}
and the numeral \lean{(2 : K)} elaborates as \lean{BitVec.instOfNat 64 2} (the bit pattern $x$),
not the cast of the natural number two (which is $0$): \texttt{\#guard ((2 : ℕ) : K) = 0} and
\texttt{\#guard (2 : K) ≠ 0} both pass. No current leanerVM declaration adds or multiplies a
\lean{BitVec 64} meaning integers (a \code{git grep} of \code{b435631} for \code{BitVec 64},
\code{UInt64} and 64-bit literals finds only \lean{BitVec.ofNat 64}, \lean{extractLsb'},
\lean{++}, \lean{<<<}, and \lean{UInt64} in BLAKE2s, whose operations are core's), so this is a
hazard, not a defect. The companion hazard, core's \lean{BitVec} simplification procedures firing
on $\K$, is the finding \enquote{\lean{simp} misreads $\K$ arithmetic}
(\cref{f:lib-simp-misreads-k}).

\begin{define}{\lean{BF64.Ext3}: the field $\E$}
$\E = \K[y]/(y^3 + y + 1)$, built by CompPoly's generic extension framework \lean{Ext P}: the
carrier is a dense coefficient vector of length \lean{P.d}, little-endian (index $i$ is the
coefficient of $y^i$), so an element of $\E$ is the limb triple $c_0 + c_1 y + c_2 y^2$.
\lean{ext3Params} records the modulus below its implicit leading term (\texttt{lower := \#v[1, 1, 0]}
are the coefficients of $1, y, y^2$), \lean{ofBase} embeds $\K$ as the $y^0$ limb, and
\lean{gen} is $y$.
\end{define}

\begin{leancode}
noncomputable def ext3Poly : Polynomial BF64 := X ^ 3 + X + 1
...
theorem ext3Poly_irreducible : Irreducible ext3Poly :=
  Polynomial.irreducible_of_degree_le_three_of_not_isRoot
    (by rw [ext3Poly_natDegree]; decide) ext3Poly_no_root
...
def ext3Params : ExtensionParams BF64 where
  d := 3
  two_le := by norm_num
  lower := #v[1, 1, 0]
  q := 2 ^ 64
  card_eq := card_bf64
...
abbrev Ext3 : Type := Ext ext3Params
...
@[simp] theorem card_ext3 : Fintype.card Ext3 = 2 ^ 192 := by
  rw [Ext.card_ext, ext3Params_q, ext3Params_d, ← pow_mul]
\end{leancode}

\src{\file{CompPoly/}\allowbreak\file{Fields/}\allowbreak\file{Binary/}\allowbreak\file{BF64/}\allowbreak\file{Ext3.lean:56-57}, \file{:125-127}, \file{:133-140}, \file{:165-168}, \file{:198-200}, CompPoly \code{3468b38c}}

Irreducibility: a root $a$ of $y^3+y+1$ has $a^7 = 1$; $\gcd(7, 2^{64}-1) = 1$
(\lean{decide +kernel}), so $a = 1$, which is not a root (\file{Ext3.lean:78-122}).
\lean{ext3Params_poly} (\file{:148-160}) proves that the coefficient vector
\texttt{\#v[1, 1, 0]} denotes $y^3 + y + 1$.

\begin{leancode}
/--
The carrier of the extension `F[X] / f`: a dense coefficient vector of length `P.d`,
little-endian (index `i` is the coefficient of `X^i`).
-/
def Ext {F : Type*} [Field F] [Fintype F] (P : ExtensionParams F) : Type _ := Vector F P.d
...
@[inline] def ofBase (c : F) : Ext P := ofFn fun i => if (i : ℕ) = 0 then c else 0
...
def gen : Ext P := ofFn fun i => if (i : ℕ) = 1 then 1 else 0
\end{leancode}

\src{\file{CompPoly/}\allowbreak\file{Fields/}\allowbreak\file{Extension/}\allowbreak\file{Defs.lean:132-136}, \file{:185}, \file{:188}, CompPoly \code{3468b38c}}

\lean{Ext} is a plain \lean{def} (not reducible), so \lean{Vector}'s instances do not leak onto
$\E$. Its multiplication folds products by the reduced monomials $X^k \bmod f$
(\file{Defs.lean:227-233}; the compiled code uses the table version \lean{mulTbl}, swapped in by a
proved \lean{@[csimp]} equation \lean{mul_eq_mulTbl}, \file{:267-272}). The ring laws are
transported from \lean{AdjoinRoot P.poly} through a proved bijection
(\file{Extension/}\allowbreak\file{Field.lean:73-107}, \lean{ringEquivQuot}); inversion is Fermat,
$x^{-1} = x^{q^d - 2}$ (\file{Field.lean:127}); the \lean{Field} instance needs
\lean{[Fact (Irreducible P.}\allowbreak\lean{poly)]} (\file{:159-166}), supplied by \file{Ext3.lean:162-163}.
\lean{Fintype (Ext P)} is by the coefficient bijection (\file{Field.lean:54-63}),
\lean{card_}\allowbreak\lean{ext : Fintype.}\allowbreak\lean{card (Ext P) = P.}\allowbreak\lean{q ^ P.d}. The embedding $\K \to \E$ is
\lean{Ext.ofBase}, packaged as \lean{algebraMap}:

\begin{leancode}
/-- The extension is an `F`-algebra. -/
instance instAlgebra : Algebra F (Ext P) where
  algebraMap := ofBaseRingHom P
  commutes' _ _ := mul_comm _ _
  smul_def' c x :=
    toQuot_injective (by simp only [toQuot_smul, toQuot_mul, ofBaseRingHom_apply, toQuot_ofBase])

@[simp] theorem algebraMap_eq_ofBase (c : F) : algebraMap F (Ext P) c = ofBase c := rfl
\end{leancode}

\src{\file{CompPoly/}\allowbreak\file{Fields/}\allowbreak\file{Extension/}\allowbreak\file{Bridge.lean:363-370}, CompPoly \code{3468b38c}}

\lean{DecidableEq (Ext P)} and \lean{BEq (Ext P)} are \lean{Vector}'s (\file{Defs.lean:282-284});
there is no \lean{LawfulBEq (Ext P)} at the old pin (\cref{sec:gen-lib-comppoly-mv}).

\subsection{Are these leanVM's fields? Yes, bit for bit}\label{sec:gen-lib-comppoly-fidelity}

A formalization over another presentation of $\mathrm{GF}(2^{64})$ would say nothing about the
deployed verifier's arithmetic. \Cref{tab:gen-lib-fields} checks the fields on four sources.
Specification: \file{doc/}\allowbreak\file{leanvm/}\allowbreak\file{body/}\allowbreak\file{02-vm-specification.tex:6-9}; Rust: under
\file{crates/}\allowbreak\file{primitives/}\allowbreak\file{src/}\allowbreak\file{field/}; Python: \file{python-verifier/}\allowbreak\file{verifier.py}; all at leanVM's
pin \code{a386121f}.

{\footnotesize
\begin{longtable}{@{}p{0.08\linewidth}p{0.14\linewidth}p{0.26\linewidth}p{0.21\linewidth}p{0.19\linewidth}@{}}
\caption{The fields $\K$ and $\E$ in the specification, the Rust, the Python and CompPoly/leanerVM.}\label{tab:gen-lib-fields}\\
\toprule
& Specification & Rust & Python & CompPoly / leanerVM \\
\midrule
\endfirsthead
\toprule
& Specification & Rust & Python & CompPoly / leanerVM \\
\midrule
\endhead
\bottomrule
\endfoot
$\K$ modulus & $x^{64}+x^4+x^3+x+1$
  & \file{gf2_64.rs:2} \enquote{\code{K = GF(2)[x]/(x^64 + x^4 + x^3 + x + 1)}}, reduction
    \code{R64 = 0x1B} (\file{gf2_64x3.rs:30})
  & \file{:39} \enquote{\code{GF(2^64) = F2[x]/(x^64 + x^4 + x^3 + x + 1)}}, fold
    \code{high ^ high<<1 ^ high<<3 ^ high<<4} (\file{:31})
  & \lean{basePoly} (\file{BF64/}\allowbreak\file{Basic.lean:102}) \\
$\K$ encoding & ---
  & \file{gf2_64.rs:19} \enquote{\code{bit i = coefficient of x^i}}, \code{F64(pub u64)}
  & \code{K(value: int)}, 64-bit (\file{:39-46})
  & \lean{BitVec 64}, bit $i$ = coefficient of $x^i$ (\file{Impl.lean:52-54}) \\
$\K$ addition & ---
  & \code{self.0 ^ rhs.0} (\file{gf2_64.rs:75})
  & \code{self.value ^ rhs.value} (\file{:67})
  & \texttt{a \textasciicircum\textasciicircum\textasciicircum{} b} (\file{Impl.lean:62}) \\
$\E$ modulus & $\K[y]/(y^3+y+1)$
  & \file{gf2_64x3.rs:3} \enquote{\code{F192 = K[y]/(y^3 + y + 1)}}
  & \file{:92} \enquote{\code{K[y]/(y^3 + y + 1)}}
  & \lean{ext3Poly}, \texttt{lower := \#v[1, 1, 0]} \\
$\E$ limbs & ---
  & \code{F192 { c0, c1, c2 }}, \enquote{\code{c0 + c1·y + c2·y²}} (\file{gf2_64x3.rs:5,34-37})
  & \code{E(c0, c1, c2)}, bytes \code{<3Q} (\file{:94-107})
  & \lean{Vector K 3}, index $i$ = coefficient of $y^i$ \\
$y$ & ---
  & \code{F192::Y = (0, 1, 0)} (\file{gf2_64x3.rs:44})
  & \code{Y = E(0, 1)} (\file{:183})
  & \lean{Ext.gen} = limbs $(0,1,0)$ \\
$\K \hookrightarrow \E$ & ---
  & limb \code{c0}
  & \code{E(lifted)} = $(k, 0, 0)$ (\file{:111-115})
  & \lean{ofBase} = $(c, 0, 0)$ \\
generator & \enquote{a generator of order $2^{64} - 1$} (\file{:12})
  & \code{F64::G = F64(2)} (\file{gf2_64.rs:28})
  & \code{GEN = E(2)} (\file{:182})
  & \lean{g : K := 0x2} (\file{Generator.lean:259}), order proved \\
\end{longtable}}

Beyond the matching definitions, the probe \code{FieldFidelity} ran the Rust's own reference
vectors through leanerVM's compiled $\K$ and $\E$ arithmetic at the old pins, and every
\texttt{\#guard} passed (the three products of \file{gf2_64.rs:271-275}, the four products and
squares of \file{gf2_64x3.rs:991-1016}, $y \cdot y \cdot y = y + 1$, $x^{63} \cdot x =
\mathtt{0x1B}$, $0^{-1} = 0$) while a mutated product and a mutated limb were rejected
(\cref{ch:probes}). Python's \code{E.__mul__} (\file{:145-152}) folds
\code{p0+p3, p1+p3+p4, p2+p4}, which is exactly leanerVM's \lean{mul_limbs}
(\file{LeanerVM/}\allowbreak\file{Parameters/}\allowbreak\file{Field.lean:206-210}).

\textbf{Conclusion.} leanerVM's $\K$ and $\E$ are the specification's fields \emph{with the
Rust's and Python's encodings}: the identification is the identity on 64-bit words and on limb
triples. (leanVM also uses $\mathrm{GF}(2^8)$ with the AES modulus and an embedding
$\varphi_8 : \mathrm{GF}(2^8) \to \K$ inside Flock, \file{crates/}\allowbreak\file{primitives/}\allowbreak\file{src/}\allowbreak\file{field/}\allowbreak\file{gf2_8.rs},
\file{phi8_tower.rs:15-24}, \file{c-flock-protocol.tex:35}; CompPoly has no such field at the old
pin and the blueprint says so, \file{protocol-blueprint.md:281}.)

\subsection{leanerVM's field modules (at \texorpdfstring{\texttt{b435631}}{b435631})}\label{sec:gen-lib-comppoly-leanervm}

$\K$ and $\E$ are \emph{abbreviations} of CompPoly's types, so leanerVM adds no field of its own;
\lean{y} is CompPoly's generator; \lean{ofK} is \lean{algebraMap K E} (probe \code{Layer0}:
\lean{example (a : K) : ofK a = algebraMap K E a := rfl} is accepted). The file's other
declarations are limb lemmas (\lean{limb_add}, \lean{mul_limbs}, \lean{isInK_iff}, \dots), all
proved.

\begin{leancode}
/-- `K = GF(2^64)`, CompPoly's `BF64`: a `BitVec 64` with bit `i` the coefficient of `x^i`. -/
abbrev K : Type := BF64

/-- `E = K[y]/(y^3 + y + 1)`, CompPoly's `BF64.Ext3`: a `Vector K 3` in limb order
`c0 + c1·y + c2·y²`. -/
abbrev E : Type := BF64.Ext3

/-- The adjoined root `y` of `y^3 + y + 1`. -/
def y : E := BF64.ext3Gen

/-- The defining relation of `E`: `y^3 = y + 1`. -/
theorem y_pow_three : y ^ 3 = y + 1 := BF64.ext3Gen_pow_three
...
/-- The embedding `K ↪ E` as the `y^0` limb; it is `algebraMap K E` by `rfl`. -/
abbrev ofK (a : K) : E := Ext.ofBase a
...
/-- The word `c0 + c1·y + c2·y²`, as the literal limb vector. -/
def E.ofLimbs (c0 c1 c2 : K) : E := Ext.ofVector #v[c0, c1, c2]
\end{leancode}

\src{\file{LeanerVM/}\allowbreak\file{Parameters/}\allowbreak\file{Field.lean:70-92}, leanerVM \code{b435631}}

\file{LeanerVM/}\allowbreak\file{Parameters/}\allowbreak\file{Generator.lean:258-311}: \lean{g : K := 0x2}, and
\lean{orderOf_}\allowbreak\lean{g : orderOf g = 2^64 − 1} by Lagrange plus seven \lean{decide +kernel} checks
$g^{(2^{64}-1)/p} \neq 1$ over the prime factors $3 \cdot 5 \cdot 17 \cdot 257 \cdot 641 \cdot
65537 \cdot 6700417$. The generator is certified, not assumed.

\file{LeanerVM/}\allowbreak\file{Parameters/}\allowbreak\file{CleanField.lean:360-372} supplies Clean's field interface for $\K$, so
Clean's $\K$ has CompPoly's field structure (one field structure, by \lean{rfl}). Clean's
\lean{val}/\lean{fromNat} read $\K$ as the integers $0, \dots, 2^{64}-1$, which is \emph{not} a
ring homomorphism in characteristic two; the file's docstring records that Clean's core never
consumes them.

\begin{leancode}
instance instFiniteFieldK : FiniteField K where
  val := BitVec.toNat
  fromNat n := BitVec.ofNat 64 n
  size := 2 ^ 64
  ...
theorem instFiniteFieldK_toField : (instFiniteFieldK.toField : Field K) = BF64.instField := rfl
\end{leancode}

\src{\file{LeanerVM/}\allowbreak\file{Parameters/}\allowbreak\file{CleanField.lean:360-372}, leanerVM \code{b435631}}

\subsection{At the new pin \texorpdfstring{\texttt{572f9973}}{572f9973}}\label{sec:gen-lib-comppoly-new}

CompPoly moved by 64 commits, from \code{3468b38c} to \code{572f9973}. What changes for the proof
system:
\begin{itemize}
\item \textbf{\lean{BF64} is a structure} (\#329, \enquote{isolate the BF64 presentation},
  \code{e5f87c8}): \lean{structure BF64 where toBitVec : BitVec 64 deriving DecidableEq, BEq}
  (\file{CompPoly/}\allowbreak\file{Fields/}\allowbreak\file{Binary/}\allowbreak\file{BF64/}\allowbreak\file{Impl.lean:57-61} at \code{572f9973}), with \lean{ofBitVec}.
  CompPoly's field instances no longer leak onto \lean{BitVec 64}; core's \lean{BitVec}
  simplification procedures no longer see $\K$; and \lean{FinEnum (BitVec 64)} no longer gives a
  sampler on $\K$ (\cref{sec:gen-l0-samplers}).
\item \textbf{Numerals change meaning.} There is no \lean{OfNat BF64} instance; a numeral
  $n \geq 2$ elaborates through Mathlib's \lean{NatCast} numeral instance, and the cast is by
  parity (\file{Impl.lean:223-227} at \code{572f9973}). So $(2 : \K) = 0$, $(3 : \K) = 1$,
  $(\mathtt{0x1B} : \K) = 1$: the characteristic-two meaning. Bit patterns are written
  \lean{K.ofBits n}, which leanerVM defines at \code{144c5aa}
  (\file{LeanerVM/}\allowbreak\file{Parameters/}\allowbreak\file{Field.lean}).
\end{itemize}

\begin{leancode}
/-- Cast a natural number by its parity in characteristic two. -/
@[inline] def natCast (n : ℕ) : BF64 := if n % 2 = 0 then 0 else 1

instance : NatCast BF64 := ⟨natCast⟩
\end{leancode}

\src{\file{CompPoly/}\allowbreak\file{Fields/}\allowbreak\file{Binary/}\allowbreak\file{BF64/}\allowbreak\file{Impl.lean:223-227}, CompPoly \code{572f9973}}


\begin{leancode}
def K.ofBits (n : ℕ) : K := BF64.ofBitVec (BitVec.ofNat 64 n)
\end{leancode}

\src{\file{LeanerVM/}\allowbreak\file{Parameters/}\allowbreak\file{Field.lean}, leanerVM \code{144c5aa}}

\paragraph{Effect on probes and fixtures.} Every numeral $\geq 2$ in $\K$ written at the old pin
\emph{compiles at the new pin with another value}. The probe \code{FieldFidelity} at the new pin
would read \lean{(0x01090913877ed8ed : K) * 0x66ab35ac2768468f = 0x50c4519dc383744a} as
$1 \cdot 1 = 0$ (the parities of the three words) and fail; its \texttt{\#guard (2 : K) ≠ 0} would
fail, \texttt{\#guard ((2 : ℕ) : K) = 0} would pass, \lean{g = (2 : K)} would fail (\lean{g} is
\lean{K.ofBits 0x2} at \code{144c5aa}). Its conclusions are about the old pin and stand there;
at the new pin the field and its encoding are the same (\lean{ofBitVec} of the same words), so
re-running it with \lean{K.ofBits} on every literal would be the check (unverified: no run is
possible now). The probe \code{NumeralHazard} is moot at the new pin. Pull request \#61 rewrote
leanerVM's fixtures with \lean{K.ofBits} (for example \file{tests/}\allowbreak\file{LeanerVMTests/}\allowbreak\file{Protocol/}\allowbreak\file{Field.lean}
at \code{144c5aa}: \texttt{⟨\#v[1, K.ofBits 2, K.ofBits 3, K.ofBits 4]⟩}) and pins the new meaning
(\file{tests/}\allowbreak\file{LeanerVMTests/}\allowbreak\file{Parameters/}\allowbreak\file{Field.lean:50-51} at \code{144c5aa}: \texttt{\#guard (2 : K) = 0},
\texttt{\#guard K.ofBits 2 ≠ (2 : K)}). A scan of \code{144c5aa} for \lean{ofK n}, \lean{(n : K)},
\lean{(n : E)}, \lean{: K := n}, \lean{E.ofLimbs n} with $n \geq 2$ or a hex literal finds only
those two intended lines and \lean{E.ofLimbs 0x0000000000000000 …} ($= 0$) in
\file{tests/}\allowbreak\file{LeanerVMTests/}\allowbreak\file{Semantics/}\allowbreak\file{Blake2s.lean:197}; numerals inside vectors and other shapes
were not scanned (unverified).

\paragraph{The other changes.}
\begin{itemize}
\item \#330, \enquote{separate extension arithmetic from finite certificates} (\code{cf340b2}):
  \lean{ExtensionParams F} no longer takes \lean{[Fintype F]}; its \lean{q} is uncertified data
  (\enquote{Raw arithmetic does not certify this value. Field laws require a separate proof that
  \lean{Nat.card F = q}}, \file{CompPoly/}\allowbreak\file{Fields/}\allowbreak\file{Extension/}\allowbreak\file{Arithmetic.lean:57-61} at
  \code{572f9973}), supplied for \lean{Ext3} by
  \lean{instance : Fact (Nat.}\allowbreak\lean{card BF64 = ext3Params.}\allowbreak\lean{q) := ⟨nat_}\allowbreak\lean{card_}\allowbreak\lean{bf64⟩}
  (\file{Ext3.lean:143-144}).
\item \#331, \enquote{avoid eager BF64 enumeration at native startup} (\code{290c351}): no
  \lean{Fintype BF64}; there is \lean{instance : Finite BF64} and
  \lean{theorem card_}\allowbreak\lean{bf64 [Fintype BF64] : Fintype.}\allowbreak\lean{card BF64 = 2 ^ 64} (\file{Impl.lean:459-468}),
  likewise \lean{card_ext3 [Fintype Ext3]} (\file{Ext3.lean:206-208}). leanerVM supplies, at
  \code{144c5aa}, \lean{noncomputable instance : Fintype K := Fintype.}\allowbreak\lean{ofFinite K} (\enquote{A
  proof-only enumeration of $\K$; executable code uses its explicit bit coordinates}), from which
  CompPoly's \lean{Ext.}\allowbreak\lean{instFintype [Fintype F]} (\file{Extension/}\allowbreak\file{Cardinality.lean:52}) gives
  \lean{Fintype E}. \textbf{So \lean{Fintype K} is now proof-only}: \lean{Fintype.card K} in
  statements is unchanged in meaning (a subsingleton class), no executable can evaluate it, and
  the startup out-of-memory failure of the eager instance is gone. The blocker on executables
  (leanISA status finding P3) is lifted by the new pin; the blueprint's executable \lean{verify}
  and the differential tests are unblocked on this count (other blockers not assessed).
\item leanerVM also adds, at \code{144c5aa}, its own \lean{instDecidableEqK} (decided on
  \lean{toBitVec}) beside CompPoly's derived \lean{DecidableEq BF64}: two instances, equal as
  propositions (\lean{Decidable} is a subsingleton), chosen for kernel reduction per its
  docstring. Harmless; one more trusted-surface line.
\item \lean{LawfulBEq (Ext P)} now exists (\file{Extension/}\allowbreak\file{Arithmetic.lean:380}, given
  \lean{LawfulBEq F}), so \lean{CMvPolynomial n E} gets \lean{C}, \lean{X} and ring operations at
  the new pin; the blueprint's \enquote{no \lean{LawfulBEq E} at the pin}
  (\file{protocol-blueprint.md:570}) is stale after \#61 (not re-checked by instance search: no
  run).
\item New fields: the AES polynomial field $\mathrm{GF}(2^8)$, modulus $X^8 + X^4 + X^3 + X + 1$
  (\file{CompPoly/}\allowbreak\file{Fields/}\allowbreak\file{Binary/}\allowbreak\file{Aes/}\allowbreak\file{Basic.lean}, \#356), and its embedding into GHASH
  $\mathrm{GF}(2^{128})$ (\#357). leanVM's Flock uses the same $\mathrm{GF}(2^8)$ and an embedding
  $\varphi_8 : \mathrm{GF}(2^8) \to \K$; CompPoly has the source field now but no embedding into
  \lean{BF64}, so the blueprint's gap statement (\file{protocol-blueprint.md:281}, \enquote{no
  \lean{GF(2^8) ↪ E} embedding}) remains true but should say the field itself has landed.
\item Serialization (\#373--\#376: byte codecs for every field carrier, \enquote{exact fiber
  counts and the bias bound of the challenge decoder}): relevant to the compiled verifier's
  transcript (Layer 12) and to modelling \code{sample() -> F192} from a hash output; not assessed
  further.
\end{itemize}

\paragraph{leanerVM's CompPoly-bound modules and upstream.} The four modules
\file{Multilinear.lean}, \file{BitProductTable.lean}, \file{Stacking.lean} and
\file{AmbientStacking.lean} under \file{LeanerVM/}\allowbreak\file{Protocol/}\allowbreak\file{ToCompPoly/} were checked by name and by content keyword (\lean{sumCube}, hypercube sums,
\lean{hadamard}/pointwise products, \lean{evalMle_eq_sum}, \lean{evalMle_split}/append,
\lean{bitProductTable}/\lean{powersTable}, stacking) in CompPoly \code{3468b38c} and
\code{572f9973} (\file{CompPoly/}\allowbreak\file{Multilinear/}) and ArkLib \code{dca90385} and \code{7653a901}
(\file{ArkLib/}\allowbreak\file{ToCompPoly/}): \textbf{no counterpart} of any of their 80-odd declarations. What
CompPoly added between the pins in \file{Multilinear/}\allowbreak\file{Basic.lean} is linearity (\lean{eval_add},
\lean{eval_zero}, \lean{eval_smul} for both representations, \lean{map_eval},
\lean{evalWithProducts}), which the local modules do not restate. The local modules remain
CompPoly candidates. One name clash is inferred, not run: CompPoly \code{572f9973} now declares
\lean{CompPoly.}\allowbreak\lean{CMlPolynomialEval.}\allowbreak\lean{eval_}\allowbreak\lean{zero} (\file{CompPoly/}\allowbreak\file{Multilinear/}\allowbreak\file{Basic.lean:719}), and
ArkLib \code{7653a901} declares the same name in \file{ArkLib/}\allowbreak\file{ToCompPoly/}\allowbreak\file{Multilinear/}\allowbreak\file{Basic.lean:35-45};
ArkLib builds against its own CompPoly pin \code{df591bb8}, which lacks it, and leanerVM's root
pin \code{572f9973} wins the resolution. So a leanerVM file importing
\lean{ArkLib.}\allowbreak\lean{ToCompPoly.}\allowbreak\lean{Multilinear.}\allowbreak\lean{Basic} would fail with a duplicate declaration. leanerVM
imports none of these modules today, but the blueprint names that module's
\lean{eval_eq_MvPolynomial_MLE} as \emph{the} bridge from tables to Mathlib multilinears
(\file{protocol-blueprint.md:246}); see \cref{f:lib-mle-bridge-unimportable}.

\par\endgroup